\documentclass[12pt]{article}
\usepackage[utf8]{inputenc}
\usepackage{bbm}
\usepackage{geometry}
\geometry{a4paper,scale=0.75}
\linespread{1.5}
\usepackage{graphicx} 
\usepackage{float} 
\usepackage{subcaption} 
\usepackage{enumerate}
\usepackage{amsmath}
\usepackage{array}
\newcolumntype{P}[1]{>{\raggedright\arraybackslash}p{#1}}
\usepackage{booktabs}
\usepackage{multirow}
\usepackage{amsfonts}
\usepackage[english]{babel}
\usepackage{amsthm}
\usepackage{dcolumn}
\usepackage{multicol}
\usepackage{stfloats}
\usepackage{placeins}
\usepackage{lscape}
\usepackage[figuresright]{rotating}
\RequirePackage{pdflscape}
\usepackage[toc,page]{appendix}
\usepackage{geometry}
\usepackage{longtable}
\usepackage{comment}
\usepackage{xcolor}
% \usepackage{amsmath}
\usepackage{amssymb}
\usepackage{empheq}
\usepackage{newtxtext,newtxmath}

% -------- enumerated sub-labels (a), (b), … --
\usepackage{enumitem}
\setlist[enumerate,1]{label=(\alph*),ref=\alph*}
% ---------------------------------------------

\usepackage{hyperref}
\hypersetup{hidelinks,
	colorlinks=true,
	allcolors=black,
	pdfstartview=Fit,
	breaklinks=true}
\usepackage{csquotes}
\usepackage{natbib}
\newtheorem{definition}{Definition}
\newtheorem{theorem}{Theorem}
%\newtheorem{proposition}[theorem]{Proposition}
\newtheorem{proposition}{Proposition}[section]
%\newtheorem{lemma}[theorem]{Lemma}
\newtheorem{lemma}{Lemma}[section]
%\newtheorem{corollary}[theorem]{Corollary}
\newtheorem{corollary}{Corollary}[section]
\newtheorem*{remark}{Remark}
\newtheorem{example}{Example}
\newtheorem{exercise}{Exercise}[section]
\numberwithin{equation}{section}
\newtheorem{assumption}{Assumption}[section] % number within sections

\title{Production Networks and Formation of Aggregate Uncertainty}
\author{
Xinyu Wei\footnotemark[1]\\\texttt{TBA}
\and
Harlly Zhou\footnotemark[2]\\\texttt{hzhouci@ust.hk}
}
\date{\today}

\begin{document}

\begin{titlepage}
    \begin{center}
        {\LARGE\bfseries Credit Cycles in Production Networks: How Uncertainty Shocks Travel\par}
        \vspace{1.5em}
        {\large
        Xinyu Wei\footnotemark[1] (\texttt{TBA})\\
        Harlly Zhou\footnotemark[2] (\texttt{hzhouci@ust.hk})\\
        }
        \vspace{1em}
        \today
    \end{center}
    \footnotetext[1]{Department of Economics, University of Rochester.}
    \footnotetext[2]{Department of Economics, Hong Kong University of Science and Technology.}
    \vspace{2em}
    \begin{abstract}
        \sloppy
        \noindent 
         \vspace{1em}

        \noindent\textbf{JEL Codes:} 


        \noindent\textbf{Keywords:} 
        \fussy
    \end{abstract}
    \thispagestyle{empty}
    \setcounter{page}{0}
\end{titlepage}


\section{Benchmark Model}\label{sec:benchmark_model}
\subsection{Model Setup}
\subsubsection{Sectors and Firms}

Consider a closed economy with $N$ industries, indexed by $i=1,\cdots,N$. Each industry has a representative firm that produces a differentiated good $y_{i,t}$ using capital $k_{i,t}$, labor $n_{i,t}$, and intermediate inputs from other industries, $\{x_{ij,t}\}_{j=1}^{N}$. The production function is
\begin{align}\label{eq:prod_func}
    y_{i,t}
    =
    z_{i,t}\zeta_i
    \left(
        k_{i,t}^{\eta_i}n_{i,t}^{1-\eta_i}
    \right)^{a_i}
    \left[
        \prod_{j=1}^{N}
        x_{ij,t}^{\alpha_{ij}}
    \right]^{1-a_i},
\end{align}
where $z_{i,t}$ is the productivity of industry $i$, $\zeta_i$ is a normalizing factor, and $x_{ij,t}$ is the quantity of good $j$ used by industry $i$ as an intermediate input. Assume that $\sum_{j=1}^{N}\alpha_{ij}=1$.

Capital evolves according to
\begin{align}\label{eq:capital_accumulation}
    k_{i,t+1}
    = (1-\delta) \underbrace{\xi_{i,t} k_{i,t}}_{\equiv \tilde{k}_{i,t+1}} + i_{i,t+1},
\end{align}
where $\delta$ is the physical capital depreciation rate, $i_{i,t+1}$ is investment, and $\xi_{i,t}\in[0,1]$ is the proportion of physical capital inherited from the previous period after liquidation. Investment is chosen before productivity is realized. In the benchmark model, investment is not subject to adjustment costs.

Productivity follows an $AR(1)$ process:
\begin{align}\label{eq:productivity_process}
    \log z_{i,t}
    = \rho^z_{i} \log z_{i,t-1}
    +\sigma_{i,t}\epsilon^z_{i,t},
\end{align}
where $\rho^z_{i}\in(0,1)$ is the persistence of productivity, $\sigma_{i,t}^2$ is the innovation variance, and $\epsilon^z_{i,t}$ is the productivity innovation. Let $\boldsymbol \sigma_t = (\sigma_{1,t},\ldots,\sigma_{N,t})'\in \mathbb{R}^N$. The innovation variance in each sector evolves according to
\begin{align}\label{eq:innovation_variance_process}
    \log \sigma_{i,t} = \rho^\sigma_{i} \log \sigma_{i,t-1} + \epsilon^\sigma_{i,t},
\end{align}
where $\rho^\sigma_{i}\in(0,1)$ is the persistence of the innovation variance. Let
\begin{align*}
    \boldsymbol\epsilon^z_t
    &=
    (\epsilon^z_{1,t},\ldots,\epsilon^z_{N,t})'\in \mathbb{R}^N,
    \qquad
    \boldsymbol\epsilon^z_t
    \sim
    \mathcal N(\mathbf 0, \mathbf{I_N}), \\
    \boldsymbol\epsilon^\sigma_t
    &=
    (\epsilon^\sigma_{1,t},\ldots,\epsilon^\sigma_{N,t})'\in \mathbb{R}^N,
    \qquad
    \boldsymbol\epsilon^\sigma_t
    \sim
    \mathcal N(\mathbf 0, \mathbf{\Sigma^\sigma}),
\end{align*}
where $\mathbf{I_N}$ is the $N\times N$ identity matrix and $\mathbf{\Sigma^\sigma}$ is the $N\times N$ variance-covariance matrix. For each $i$, $\epsilon^z_{i,t}$ and $\epsilon^\sigma_{i,t}$ have mean zero and unit variance. The uncertainty shock $\epsilon^\sigma_{i,t}$ is realized with the innovation variance $\sigma_{i,t}$, while the productivity shock $\epsilon^z_{i,t}$ is realized with productivity $z_{i,t}$.



At the beginning of each period $t$, the innovation variance is realized, but the productivity innovation and productivity are not yet realized. At this stage, firm $i$ chooses borrowing $b^f_{i,t}$ from households to maximize its lifetime value. Borrowing finances labor costs, investment costs, and intermediate input costs. The firm inherits depreciated and post-liquidation physical capital from the previous period, $(1-\delta)\tilde{k}_{i,t}$, and simultaneously chooses investment $i_{i,t}$ so that $k_{i,t} = (1-\delta) \tilde{k}_{i,t} + i_{i,t}$.

After productivity is realized, prices adjust flexibly to clear goods markets. Given the borrowing and investment chosen in the first stage, firms choose labor demand $n_{i,t}$ and intermediate inputs $x_{i1,t},\ldots,x_{iN,t}$.

Goods-market clearing in industry $i$ is
\begin{align}\label{eq:goods_market_clearing}
    y_{i,t}
    =
    c_{i,t}+\sum_{j=1}^{N}x_{ji,t}.
\end{align}
After sales are realized, firm $i$ earns nominal revenue $p_{i,t} y_{i,t}$, subject to a proportional revenue tax $\tau_{i,t}$. Revenue is observable, and creditors have seniority over revenue and pledgeable assets.

The firm repays borrowing $b^f_{i,t}$ at the interest rate $r_{i,t}$, where $r_{i,t}$ is the sector-specific interest rate that clears the debt market. Default occurs when the firm lacks sufficient enforceable resources. Three cases can arise:
\begin{enumerate}
    \item The firm's revenue is sufficient for repayment. The firm does not default and rebates the remaining profit to households. The set of productivity realizations that generates this case is $\Xi^1_{i,t}$:
    \begin{align*}
        \Xi^1_{i,t} &:= \left\{z_{i,t}: (1-\tau_{i,t})p_{i,t} y_{i,t} \geq (1+r_{i,t})b^f_{i,t}\right\}.
    \end{align*}
    \item The firm's revenue is insufficient for repayment, so the firm liquidates part of its physical capital $k_{i,t}$. If the liquidated capital is enough to repay, the firm earns zero profit and continues production with the remaining capital, $\xi_{i,t} k_{i,t}$. The set of productivity realizations that generates this case is $\Xi^2_{i,t}$:
    \begin{align*}
        \Xi^2_{i,t} &:= \left\{z_{i,t}: (1+r_{i,t})b^f_{i,t} > (1-\tau_{i,t})p_{i,t} y_{i,t} \geq (1+r_{i,t})b^f_{i,t} - p_{K,t} k_{i,t}\right\}.
    \end{align*}
    \item The firm's revenue is insufficient for repayment, and full liquidation of capital is still not enough to repay. The firm then continues production with zero inherited capital. In the next period, it must use working capital for investment and capital accumulation. The set of productivity realizations that generates this case is $\Xi^3_{i,t}$:
    \begin{align*}
        \Xi^3_{i,t} &:= \left\{z_{i,t}: (1-\tau_{i,t})p_{i,t} y_{i,t} + p_{K,t} k_{i,t} < (1+r_{i,t})b^f_{i,t}\right\}.
    \end{align*}
\end{enumerate}

In summary, firms make decisions in two stages. After innovation variances are realized but before productivity is realized, firm $i$ chooses borrowing $b^f_{i,t}$ and investment $i_{i,t}$. Its first-stage problem is
\begin{align}\label{eq:firm_first_stage_decision}
    \mathcal{J}
    \left(\tilde{k}_{i,t}; \boldsymbol\sigma_t\right)
    =
    \max_{b^f_{i,t}, i_{i,t}}
    \mathbb{E}_{i,t}
    \Bigg[
            d_{i,t}
            +
            \mathcal{M}_{t,t+1}
            \mathcal{J}\left(\tilde{k}_{i,t+1}; \boldsymbol\sigma_{t+1}\right)
    \Bigg],
\end{align}
where $\tilde{k}_{i,t} \equiv \xi_{i,t-1}k_{i,t-1}$ is inherited capital from $t-1$ to $t$. The problem is subject to the law of motion for capital in \eqref{eq:capital_accumulation} and the working capital constraint $b^f_{i,t} \geq w_t n_{i,t} + \sum_{j=1}^{N} p_{j,t} x_{ij,t} + p_{K,t} i_{i,t}$. Here $\mathbb{E}_{i,t}$ denotes the expectation conditional on information available after innovation variances are realized but before productivity is realized. The variable $d_{i,t}$ denotes the firm's state-contingent maximized equity payout after productivity is realized:
\begin{align*}
    d_{i,t}
    =
    \max\left\{
        0,
        (1-\tau_{i,t})p_{i,t} y_{i,t}
        -
        (1+r_{i,t})b^f_{i,t}
    \right\},
\end{align*}
where the realized value is determined in the second stage.

In the second stage, after productivity is realized, the firm chooses inputs:
\begin{align}\label{eq:firm_second_stage_decision}
    \min_{n_{i,t},\{x_{ij,t}\}_{j=1}^{N}}
    w_t n_{i,t}
    + \sum_{j=1}^{N} p_{j,t} x_{ij,t}
\end{align}
subject to
\begin{align*}
    b^f_{i,t} &\geq w_t n_{i,t}
    + \sum_{j=1}^{N} p_{j,t} x_{ij,t}
    + p_{K,t} i_{i,t},
    \\
    y_{i,t}
    &=
    z_{i,t}\zeta_i
    \left(
        k_{i,t}^{\eta_i}n_{i,t}^{1-\eta_i}
    \right)^{a_i}
    \left[
        \prod_{j=1}^{N}
        x_{ij,t}^{\alpha_{ij}}
    \right]^{1-a_i},
\end{align*}
where $w_t$ is the nominal wage, $p_{j,t}$ is the price of good $j$, and $p_{K,t}$ is the price of physical capital. This yields the following log marginal cost:
\begin{align}\label{eq:marginal_cost}
    \text{mc}_{i,t} = \frac{1}{1-a_i\eta_i} \left[(1-\eta_i)a_i \text{w}_t + (1-a_i)\sum_{j=1}^N \alpha_{ij} \text{p}_{j,t} - \eta_i a_i \text{k}_{i,t} - \text{z}_{i,t}\right],
\end{align}
where the normalization factor is set as
\begin{align*}
    \zeta_i = \dfrac{[1-\eta_i a_i]^{1-\eta_i a_i}}{[(1-\eta_i)a_i]^{(1-\eta_i)a_i}\left[\prod_{j=1}^N \alpha_{ij}^{\alpha_{ij}}\right]^{1-a_i}}.
\end{align*}

For the first-stage problem, the firm's FOCs are
\begin{align*}
    [b^f_{i,t}]: \quad &\mathbb{E}_{i,t}\left[\lambda^{f}_{i,t} - (1+r_{i,t})\mathbbm{1}_{\Xi^1_{i,t}} - \mathcal{M}_{t,t+1} \frac{1+r_{i,t}}{p_{K,t}} \mathcal{J}'(\tilde{k}_{i,t+1})\mathbbm{1}_{\Xi^2_{i,t}} \right] = 0, \\
    [i_{i,t}]: \quad &\mathbb{E}_{i,t}\bigg[\lambda^{f}_{i,t}p_{K,t} \\
    &\qquad - \mathcal{M}_{t,t+1} \mathcal{J}'(\tilde{k}_{i, t+1})\left(\mathbbm{1}_{\Xi^1_{i,t} \cup \Xi^2_{i,t}} + \eta_i a_i(1-\tau_{i,t+1})p_{i,t+1}\frac{y_{i,t+1}/k_{i,t+1}}{p_{K,t+1}}\mathbbm{1}_{\Xi^2_{i,t}}\right)\bigg] = 0,
\end{align*}
where $\mathbbm{1}$ denotes the indicator function for the sets $\Xi^1_{i,t}, \Xi^2_{i,t}, \Xi^3_{i,t}$. 

The envelope condition is
\begin{align*}
    \mathcal{J}'(\tilde{k}_{i,t}) = (1-\delta) &\mathbb{E}_{i,t}\bigg[\eta_i a_i (1-\tau_{i,t})p_{i,t}\frac{y_{i,t}}{k_{i,t}}\mathbbm{1}_{\Xi^1_{i,t}} \\
    &+ \mathcal{M}_{t,t+1} \mathcal{J}'(\tilde{k}_{i, t+1})\left(\mathbbm{1}_{\Xi^1_{i,t} \cup \Xi^2_{i,t}} + \eta_i a_i(1-\tau_{i,t})p_{i,t}\frac{y_{i,t}/k_{i,t}}{p_{K,t}}\mathbbm{1}_{\Xi^2_{i,t}}\right)\bigg].
\end{align*}
Equivalently, let $f_{i,t}(z)$ denote the conditional density of $z_{i,t}$ given the information available at the first stage. Since $\log z_{i,t}=\rho_i^z\log z_{i,t-1}+\sigma_{i,t}\epsilon^z_{i,t}$ and $\epsilon^z_{i,t}\sim\mathcal N(0,1)$,
\begin{align*}
    f_{i,t}(z)
    =
    \frac{1}{z\sigma_{i,t}}
    \phi\left(
        \frac{\log z-\rho_i^z\log z_{i,t-1}}{\sigma_{i,t}}
    \right),
\end{align*}
where $\phi(\cdot)$ is the standard normal density. Define the conditional probabilities
\begin{align*}
    \pi^m_{i,t}
    :=
    \Pr_{i,t}\left(z_{i,t}\in\Xi^m_{i,t}\right)
    =
    \int_{\Xi^m_{i,t}} f_{i,t}(z)\,dz,
    \qquad
    m=1,2,3,
\end{align*}
and $\pi^{12}_{i,t}:=\pi^1_{i,t}+\pi^2_{i,t}$. Let $\mathbb{E}^m_{i,t}$ denote the expectation conditional on state $\Xi^m_{i,t}$. The first-order conditions can then be written explicitly as
\begin{align}
    [b^f_{i,t}]: \quad
    &\mathbb{E}_{i,t}\left[\lambda^f_{i,t}\right]
    -
    (1+r_{i,t})\pi^1_{i,t}
    -
    \frac{1+r_{i,t}}{p_{K,t}}
    \pi^2_{i,t}
    \mathbb{E}^{2}_{i,t}\left[
        \mathcal M_{t,t+1}
        \mathcal{J}'(\tilde{k}_{i,t+1})
    \right]
    =
    0, \label{eq:first_stage_foc_b}
    \\
    [i_{i,t}]: \quad
    &p_{K,t}\mathbb{E}_{i,t}\left[\lambda^f_{i,t}\right]
    -
    \pi^{12}_{i,t}
    \mathbb{E}^{12}_{i,t}\left[
        \mathcal M_{t,t+1}
        \mathcal{J}'(\tilde{k}_{i,t+1})
    \right]\notag \\
    &\qquad \qquad -
    \pi^{2}_{i,t}
    \mathbb{E}^2_{i,t} \left[
        \mathcal{M}_{t,t+1}\mathcal{J}'(\tilde{k}_{i,t+1})
        \frac{\partial [(1-\tau_{i,t+1})p_{i,t+1} y_{i,t+1}]/\partial k_{i,t+1}
        }{
            p_{K,t}
        }
    \right] = 0. \label{eq:first_stage_foc_i}
\end{align}
The two first-order conditions have intuitive interpretations. In the borrowing condition, $\mathbb{E}_{i,t}\left[\lambda^f_{i,t}\right]$ is the marginal value of relaxing the working capital constraint by one unit, or equivalently of raising borrowing by one unit. In state $\Xi^1_{i,t}$, the firm only repays the promised debt. In state $\Xi^2_{i,t}$, the firm must also account for the marginal value loss associated with liquidation. In state $\Xi^3_{i,t}$, there is no marginal liquidation cost. For the investment condition, consider reducing investment by one unit, which frees $p_{K,t}$ units of borrowing capacity. Investment carried into the next period affects continuation value in states $\Xi^1_{i,t}$ and $\Xi^2_{i,t}$. In state $\Xi^3_{i,t}$, all investment is liquidated, so continuation value is unaffected. In state $\Xi^2_{i,t}$, the reduction in investment also affects next-period profitability through the amount of capital carried forward.

Combining the first-order condition for $i_{i,t}$ with the envelope condition gives
\begin{align}\label{eq:first_stage_envelope}
    \mathcal{J}'(\tilde{k}_{i,t}) = (1-\delta) \left[\pi^1_{i,t}\mathbb{E}^1_{i,t} \left[\frac{\partial [(1-\tau_{i,t})p_{i,t} y_{i,t}]}{\partial k_{i,t}}\right] + \mathbb{E}_{i,t} \left[\lambda^f_{i,t} p_{K,t}\right]\right].
\end{align}
The left-hand side is the marginal value of one additional unit of physical capital carried from period $t-1$, net of depreciation. Additional physical capital relaxes the working capital constraint by $(1-\delta)p_{K,t}$ units of investment. In state $\Xi^1_{i,t}$, it also raises firm profits.

\subsubsection{Households}

There is a unit mass of representative households that supply labor and credit to firms and choose consumption quantities $\{c_{i,t}\}_{i=1}^{N}$. Household utility is
\begin{align}\label{eq:hh_utility}
    U(C_t,N_t)
    =
    \log C_t
    -
    \chi\frac{N_t^{1+\varphi}}{1+\varphi},
\end{align}
where consumption aggregates sectoral goods according to
\begin{align*}
    C_t
    :=
    \prod_{i=1}^{N}
    \left(
        \frac{c_{i,t}}{\theta_i}
    \right)^{\theta_i},
\end{align*}
with $\theta_i>0$ and $\sum_{i=1}^N \theta_i=1$. Total labor supply is $N_t=\sum_{i=1}^{N}n_{i,t}$.

The household budget constraint is
\begin{align}\label{eq:hh_budget_constraint}
    &w_tN_t
    +\sum_{i=1}^{N}
        (1+\tilde{r}_{i,t})b^h_{i,t}
    +\sum_{i=1}^N e_{i,t} (d_{i,t} + q_{i,t}) + T_t \notag\\
    &\qquad \qquad\geq
    \sum_{i=1}^{N}
        p_{i,t}c_{i,t}
    +\sum_{i=1}^{N}b^h_{i,t+1} +\sum_{i=1}^N e_{i,t+1}  q_{i,t}.
\end{align}
Here $b^h_{i,t}$ denotes the household's holdings of firm $i$'s debt, $e_{i,t} \in [0,1]$ is the share of firm $i$'s equity held by households, $q_{i,t}$ is the price of firm $i$'s equity, and $\tilde{r}_{i,t}$ is the realized return on firm $i$'s debt. In equilibrium, $e_{i,t}=1$ for all $i$ and $t$. The realized return is
\begin{align}\label{eq:hh_realized_return}
    1+\tilde{r}_{i,t} \equiv
    \min\Bigg\{
        (1+r_{i,t}),
        \dfrac{(1-\tau_{i,t})p_{i,t} y_{i,t}
        +
        p_{K,t}
        \left[
            (1-\delta)\tilde{k}_{i,t}+i_{i,t}
        \right]}{b_{i,t}}
    \Bigg\},
\end{align}
where $b^f_{i,t}$ is firm $i$'s debt issuance. In equilibrium, the market for firm debt clears:
\begin{align}\label{eq:firm_debt_market_clearing}
    b^h_{i,t}=b^f_{i,t} \equiv b_{i,t}.
\end{align}
Here $D_t:= \sum_{i=1}^N d_{i,t}$ is total nominal firm profit, and $T_t$ is a lump-sum government transfer. Total consumption expenditure is
\begin{align*}
    P_tC_t
    :=
    \sum_{i=1}^{N}
        p_{i,t}c_{i,t},
\end{align*}
where $P_t$ is the aggregate price index. Following \cite{herskovicNetworksProductionAsset2018}, we normalize prices so that the consumption aggregator is the numeraire, that is, $P_t=1$.

In summary, the household problem is
\begin{align}\label{eq:hh_problem}
    \max_{\left\{\left\{c_{i,t}, b^h_{i,t+1}, e_{i,t+1}\right\}_{i=1}^{N}, N_t\right\}} \mathbb{E}_s \sum_{t=s}^{+\infty} \beta^{t-s} U(C_t, N_t)
\end{align}
subject to the budget constraint \eqref{eq:hh_budget_constraint}.

As in the standard literature, the first-order condition for $c_{i,t}$ gives consumption demand and the price index:
\begin{align*}
    c_{i,t}
    =
    \theta_i\frac{C_t}{p_{i,t}},
    \qquad \qquad
    1
    =
    \prod_{i=1}^N p_{i,t}^{\theta_i}.
\end{align*}
Labor supply follows from the first-order condition for $N_t$:
\begin{align*}
    w_t = \chi C_t N_t^{\varphi}.
\end{align*}
Forward substitution in the first-order condition for $e_{i,t+1}$ gives the equity pricing equation:
\begin{align*}
    \frac{q_{i,t}}{C_t} = \mathbb{E}_t \sum_{j=1}^{+\infty} \beta^j \frac{d_{i,t+j}}{C_{t+j}}.
\end{align*}
The stochastic discount factor therefore has the following analytical form:
\begin{align*}
    \mathcal{M}_{t, t+j} = \beta^j \frac{C_t}{C_{t+j}}.
\end{align*}

The credit-market object is the promised rate on risky firm debt. Firms choose debt issuance, households choose debt holdings, and the promised interest rate $r_{i,t+1}$ is pinned down by the risky-debt pricing condition derived from the household's first-order condition for $b^h_{i,t+1}$:
\begin{align}\label{eq:risky_debt_pricing}
    1
    =
    \mathbb{E}_{i,t}
    \left[
        \mathcal M_{t,t+1}
        (1+\tilde r_{i,t+1})
    \right].
\end{align}
Thus households do not choose firms' borrowing directly. Instead, the spread on risky debt adjusts so that the value of state-contingent repayment equals the unit cost of lending.

\subsubsection{Government}
The government collects taxes from firms, or provides subsidies to them, and rebates the resulting revenue to households as a lump-sum transfer:
\begin{align*}
    T_t
    \equiv
    \sum_{i=1}^{N}
        \tau_{i,t}p_{i,t}y_{i,t}.
\end{align*}








\clearpage
\bibliographystyle{aer}
\bibliography{network_uncertainty}


\clearpage


%%%%%%%%%%%%%%%%%%%%%%%%%%
% Publication Appendix
%%%%%%%%%%%%%%%%%%%%%%%%%%
%%%%%%%%%%%%%%%%%%%%%%%%%%%%%%%%%%%%%%%%%%%%%%%%%%%%%%%%%%%%%%%%%%%%%%%%%%%%%%%%%%%%%%%
%%%%%%%%%%%%%%%%%%%%%%%%%%%%%%%%%%%%%%%%%%%%%%%%%%%%%%%%%%%%%%%%%%%%%%%%%%%%%%%%%%%%%%%
%%%%%%%%%%%%%%%%%%%%%%%%%%%%%%%%%%%%%%%%%%%%%%%%%%%%%%%%%%%%%%%%%%%%%%%%%%%%%%%%%%%%%%%
%%%%%%%%%%%%%%%%%%%%%%%%%%%%%%%%%%%%%%%%%%%%%%%%%%%%%%%%%%%%%%%%%%%%%%%%%%%%%%%%%%%%%%%
% \clearpage
\begin{appendix}
% Formatting options
\addcontentsline{toc}{section}{Appendices}
\renewcommand{\thesection}{\Roman{section}} 
\renewcommand{\thesubsection}{\Roman{subsection}}
\renewcommand{\theequation}{\thesection\arabic{equation}}
\numberwithin{equation}{section}
% \setcounter{page}{1}
\setcounter{figure}{0} \renewcommand{\thefigure}{\Roman{section}.\arabic{figure}}
\setcounter{table}{0} \renewcommand{\thetable}{\Roman{section}.\arabic{table}}
\setcounter{subsection}{0} \renewcommand{\thesubsection}{\Roman{section}.\arabic{subsection}}
\setcounter{equation}{0} \renewcommand{\theequation}{\Roman{section}.\arabic{equation}}
\setcounter{theorem}{0} \renewcommand{\thetheorem}{\Roman{section}.\arabic{theorem}}
\setcounter{definition}{0} \renewcommand{\thedefinition}{\Roman{section}.\arabic{definition}}
% \setcounter{proposition}{0} \renewcommand{\theproposition}{\Roman{section}.\arabic{proposition}}
\setcounter{assumption}{0} \renewcommand{\theassumption}{\Roman{section}.\arabic{assumption}}
% \setcounter{hyp}{0} \renewcommand{\thehyp}{\Roman{section}.\arabic{hyp}}
\setcounter{footnote}{0}
%%%%%%%%%%%%%%%%%%%%%%%%%%%%%%%%%%%%%%%%%%%%%%%%%%%%%%%%%%%%%%%%%%%%%%%%%%%%%%%%%%%%%%%
%%%%%%%%%%%%%%%%%%%%%%%%%%%%%%%%%%%%%%%%%%%%%%%%%%%%%%%%%%%%%%%%%%%%%%%%%%%%%%%%%%%%%%%
%%%%%%%%%%%%%%%%%%%%%%%%%%%%%%%%%%%%%%%%%%%%%%%%%%%%%%%%%%%%%%%%%%%%%%%%%%%%%%%%%%%%%%%
%%%%%%%%%%%%%%%%%%%%%%%%%%%%%%%%%%%%%%%%%%%%%%%%%%%%%%%%%%%%%%%%%%%%%%%%%%%%%%%%%%%%%%%



\end{appendix}

\end{document}